\documentclass[../../main.tex]{subfiles}
\begin{document}

\paragraph{【乙】}
空間の1点$O$を通る4直線で，どの3直線も同一平面上にないようなものを考える．
このとき，4直線のいずれとも$O$以外の点で交わる平面で，4つの交点が平行四辺形の頂点になるようなものが存在することを示せ．

\paragraph{【甲】}
$AB=AC$である二等辺三角形$ABC$を考える．辺$AB$の中点を$M$とし，辺$AB$を延長した直線上に点$N$を，$AN:NB=2:1$となるようにとる．このとき
\begin{align}
  \angle BCM = \angle BCN
\end{align}
となることを示せ．ただし，点$N$は辺$AB$上にはないものとする．

\end{document}
